% ID: SEL-D0047-OA-E03
% ESTADO: cotejada. Segunda pasada visual y borrador compilado revisados.
% Fuente: Selectividad/sel_2014_matematicas_aplicadas/Matematicas_CCSS_SEPT.pdf
% SHA-256: ee5bd8e23465d55ed58339e7d9adf3d58423ac2966b53d92e39a2f9b35558ab1
% Página: 1; opción A; ejercicio original 3.

Se sabe que dos alumnos de la asignatura de Matemáticas asisten a clase, de forma independiente, el primero a un 85\% de las clases y el segundo a un 35\%. Tomado al azar un día de clase, calcule la probabilidad de cada uno de los siguientes sucesos:
\begin{enumerate}[label=\alph*)]
\item \textbf{(0.75 puntos)} Que los dos hayan asistido a clase ese día.
\item \textbf{(0.75 puntos)} Que alguno de ellos haya asistido a clase ese día.
\item \textbf{(0.5 puntos)} Que ninguno haya asistido a clase ese día.
\item \textbf{(0.5 puntos)} Que haya asistido a clase el segundo, sabiendo que el primero no ha asistido.
\end{enumerate}
